\chapter{2022年全国甲卷（文）}

\section{单选题}

\begin{problem}
已知集合 \(A=\{-2,-1,0,1,2\}\)，\(B=\{x\mid 0\le x<5\}\)，则 \(A\cap B=\)（\quad）

\choices
    {\(\{0,1,2\}\)}
    {\(\{-2,-1,0\}\)}
    {\(\{0,1\}\)}
    {\(\{1,2\}\)}

\end{problem}

\begin{answer}
A
\end{answer}

\begin{solution}
因为 \(A=\{-2,-1,0,1,2\}\)，\(B=\{0,1,2,3,4\}\)，所以
\(A\cap B=\{0,1,2\}\).
故选：\(A\).
\end{solution}

\begin{problem}
某社区通过公益讲座以普及社区居民的垃圾分类知识. 为了解讲座效果，随机抽取 \(10\) 位社区居民，让他们在讲座前和讲座后各回答一份垃圾分类知识问卷，这 \(10\) 位社区居民在讲座前和讲座后问卷答题的正确率如下图：

\bitmapfigure[max width=6.0cm,max totalheight=4.8cm]{img/2022/national_paper_a_liberal/q2_fig1.png}

则（\quad）

\choices
    {讲座前问卷答题的正确率的中位数小于 \(70\%\)}
    {讲座后问卷答题的正确率的平均数大于 \(85\%\)}
    {讲座前问卷答题的正确率的标准差小于讲座后正确率的标准差}
    {讲座后问卷答题的正确率的极差大于讲座前正确率的极差}

\end{problem}

\begin{answer}
B
\end{answer}

\begin{solution}
对于 \(A\)，讲座前问卷答题的正确率从小到大为：
\(60\%,60\%,65\%,65\%,70\%,75\%,80\%,85\%,90\%,95\%\)
所以讲座前问卷答题的正确率的中位数为
\(\dfrac{70\%+75\%}{2}=72.5\%\)
故 \(A\) 错误.

对于 \(B\)，讲座后问卷答题的正确率的平均数为：
\(\dfrac{80\%+85\%+85\%+85\%+85\%+90\%+90\%+95\%+100\%+100\%}{10}=89.5\%\)
所以 \(89.5\%>85\%\)，故 \(B\) 正确.

对于 \(C\)，由图形可知讲座前问卷答题的正确率相对分散，讲座后问卷答题的正确率相对集中，
所以讲座前问卷答题的正确率的标准差大于讲座后正确率的标准差，故 \(C\) 错误.

对于 \(D\)，讲座后问卷答题的正确率的极差为
\(100\%-80\%=20\%\)
讲座前正确率的极差为
\(95\%-60\%=35\%\)
所以讲座后问卷答题的正确率的极差小于讲座前正确率的极差，故 \(D\) 错误.
故选：\(B\).
\end{solution}

\begin{problem}
若 \(z=1+\i\)，则 \(\left|\i z+3\overline{z}\right|=\)（\quad）

\choices
    {\(4\sqrt{5}\)}
    {\(4\sqrt{2}\)}
    {\(2\sqrt{5}\)}
    {\(2\sqrt{2}\)}

\end{problem}

\begin{answer}
D
\end{answer}

\begin{solution}
因为 \(z=1+\i\)，所以 \(\overline{z}=1-\i\)，于是
\(\i z+3\overline{z} =\i\left(1+\i\right)+3\left(1-\i\right) =2-2\i\)
所以
\(\left|\i z+3\overline{z}\right| =\sqrt{2^2+\left(-2\right)^2} =2\sqrt{2}\).
故选：\(D\).
\end{solution}

\begin{problem}
如图，网格纸上绘制的是一个多面体的三视图，网格小正方形的边长为 \(1\)，则该多面体的体积为（\quad）

\bitmapfigure[max width=6.0cm,max totalheight=4.8cm]{img/2022/national_paper_a_liberal/q4_fig1.png}

\choices
    {8}
    {12}
    {16}
    {20}

\end{problem}

\begin{answer}
B
\end{answer}

\begin{solution}
由多面体的三视图得该多面体是一正四棱柱 \(ABCD-A_1B_1C_1D_1\)，四棱柱的底面是直角梯形 \(ABCD\).
由图可知 \(AB=4\)，\(AD=2\)，\(AA_1=2\)，且 \(AA_1\perp\) 平面 \(ABCD\).
所以该多面体的体积为
\(V=\dfrac{1}{2}\times\left(4+2\right)\times2\times2=12\).
故选：\(B\).
\end{solution}

\begin{problem}
将函数 \(f(x)=\sin\left(\omega x+\dfrac{\pi}{3}\right)\)（\(\omega>0\)）的图像向左平移 \(\dfrac{\pi}{2}\) 个单位长度后得到曲线 \(C\)，若 \(C\) 关于 \(y\) 轴对称，则 \(\omega\) 的最小值是（\quad）

\choices
    {\(\dfrac{1}{6}\)}
    {\(\dfrac{1}{4}\)}
    {\(\dfrac{1}{3}\)}
    {\(\dfrac{1}{2}\)}

\end{problem}

\begin{answer}
C
\end{answer}

\begin{solution}
由题意知，曲线 \(C\) 的解析式为 \(y=\sin\left[\omega\left(x+\dfrac{\pi}{2}\right)+\dfrac{\pi}{3}\right]=\sin\left(\omega x+\dfrac{\omega\pi}{2}+\dfrac{\pi}{3}\right)\).
又 \(C\) 关于 \(y\) 轴对称，所以
\(\dfrac{\omega\pi}{2}+\dfrac{\pi}{3}=\dfrac{\pi}{2}+k\pi\)，\(k\in \Z\)
解得
\(\omega=\dfrac{1}{3}+2k\)，\(k\in \Z\).
因为 \(\omega>0\)，所以当 \(k=0\) 时，\(\omega\) 的最小值为
\(\dfrac{1}{3}\).
故选：\(C\).
\end{solution}

\begin{problem}
从分别写有 \(1,2,3,4,5,6\) 的 \(6\) 张卡片中无放回随机抽取 \(2\) 张，则抽到的 \(2\) 张卡片上的数字之积是 \(4\) 的倍数的概率为（\quad）

\choices
    {\(\dfrac{1}{5}\)}
    {\(\dfrac{1}{3}\)}
    {\(\dfrac{2}{5}\)}
    {\(\dfrac{2}{3}\)}

\end{problem}

\begin{answer}
C
\end{answer}

\begin{solution}
从 \(6\) 张卡片中无放回抽取 \(2\) 张，共有 \(\left(1,2\right)\)，\(\left(1,3\right)\)，\(\left(1,4\right)\)，\(\left(1,5\right)\)，\(\left(1,6\right)\)，\(\left(2,3\right)\)，\(\left(2,4\right)\)，\(\left(2,5\right)\)，\(\left(2,6\right)\)，\(\left(3,4\right)\)，\(\left(3,5\right)\)，\(\left(3,6\right)\)，\(\left(4,5\right)\)，\(\left(4,6\right)\)，\(\left(5,6\right)\) 15 种情况.
其中数字之积为 \(4\) 的倍数的有 \(\left(1,4\right)\)，\(\left(2,4\right)\)，\(\left(2,6\right)\)，\(\left(3,4\right)\)，\(\left(4,5\right)\)，\(\left(4,6\right)\) 共 \(6\) 种情况，所以所求概率为
\(\dfrac{6}{15}=\dfrac{2}{5}\).
故选：\(C\).
\end{solution}

\begin{problem}
函数 \(y=\left(3^x-3^{-x}\right)\cos x\) 在区间 \(\left[-\dfrac{\pi}{2},\dfrac{\pi}{2}\right]\) 的图象大致为（\quad）

\choices
    {图象形态 A}
    {图象形态 B}
    {图象形态 C}
    {图象形态 D}

\end{problem}

\begin{answer}
A
\end{answer}

\begin{solution}
令
\(f(x)=\left(3^x-3^{-x}\right)\cos x\)，\(x\in \left[-\dfrac{\pi}{2},\dfrac{\pi}{2}\right]\)
则
\(f(-x)=\left(3^{-x}-3^x\right)\cos\left(-x\right)=-\left(3^x-3^{-x}\right)\cos x=-f(x)\)
所以 \(f(x)\) 为奇函数，排除 \(B,D\).
又当 \(x\in \left(0,\dfrac{\pi}{2}\right)\) 时，
\(3^x-3^{-x}>0\)，\(\cos x>0\)
所以
\(f(x)>0\)
排除 \(C\).
故选：\(A\).
\end{solution}

\begin{problem}
当 \(x=1\) 时，函数 \(f(x)=a\ln x+\dfrac{b}{x}\) 取得最大值 \(-2\)，则 \(f'(2)=\)（\quad）

\choices
    {\(-1\)}
    {\(-\dfrac{1}{2}\)}
    {\(\dfrac{1}{2}\)}
    {1}

\end{problem}

\begin{answer}
B
\end{answer}

\begin{solution}
因为函数 \(f(x)\) 在 \(x=1\) 时取得最大值 \(-2\)，所以
\(f(1)=-2\)，\(f'(1)=0\).
又
\(f(1)=a\ln1+\dfrac{b}{1}=b=-2\)
所以 \(b=-2\).
且
\(f'(x)=\dfrac{a}{x}-\dfrac{b}{x^2}\)
由 \(f'(1)=0\) 得
\(a-b=0\)
代入 \(b=-2\)，得 \(a=-2\).
所以
\(f'(x)=-\dfrac{2}{x}+\dfrac{2}{x^2}\)
从而
\(f'(2)=-1+\dfrac{1}{2}=-\dfrac{1}{2}\).
故选：\(B\).
\end{solution}

\begin{problem}
在长方体 \(ABCD-A_1B_1C_1D_1\) 中，已知 \(B_1D\) 与平面 \(ABCD\) 和平面 \(AA_1B_1B\) 所成的角均为 \(30^\circ\)，则（\quad）

\bitmapfigure[max width=6.0cm,max totalheight=4.8cm]{img/2022/national_paper_a_liberal/q9_fig1.png}

\choices
    {\(AB=2AD\)}
    {\(AB\) 与平面 \(AB_1C_1D\) 所成的角为 \(30^\circ\)}
    {\(AC=CB_1\)}
    {\(B_1D\) 与平面 \(BB_1C_1C\) 所成的角为 \(45^\circ\)}

\end{problem}

\begin{answer}
D
\end{answer}

\begin{solution}
如图所示，设 \(AB=a\)，\(AD=b\)，\(AA_1=c\).
依题意以及长方体的结构特征可知，\(B_1D\) 与平面 \(ABCD\) 所成角为 \(\angle B_1DB\)，\(B_1D\) 与平面 \(AA_1B_1B\) 所成角为 \(\angle DB_1A\).
所以
\(\sin30^\circ=\dfrac{c}{B_1D}=\dfrac{b}{B_1D}\)
故 \(b=c\).
又
\(B_1D=2c=\sqrt{a^2+b^2+c^2}\)
解得
\(a=\sqrt{2}c\).

对于 \(A\)，\(AB=a\)，\(AD=b=c\)，所以 \(AB=2AD\) 不成立，\(A\) 错误.

对于 \(B\)，过 \(B\) 作 \(BE\perp AB_1\) 于 \(E\). 因为 \(AD\perp AB\)，\(AD\perp AA_1\)，且 \(AB,AA_1\) 是平面 \(ABB_1A_1\) 内两条相交直线，所以 \(AD\perp\) 平面 \(ABB_1A_1\). 又 \(BE\) 在平面 \(ABB_1A_1\) 内，所以 \(AD\perp BE\). 由于 \(BE\perp AB_1\)，且 \(AD,AB_1\) 是平面 \(AB_1C_1D\) 内两条相交直线，所以 \(BE\perp\) 平面 \(AB_1C_1D\). 所以 \(AB\) 与平面 \(AB_1C_1D\) 所成角为 \(\angle BAE\).
因为
\(\tan\angle BAE=\dfrac{c}{a}=\dfrac{1}{\sqrt{2}}\)
所以 \(\angle BAE\ne30^\circ\)，\(B\) 错误.

对于 \(C\)，
\(AC=\sqrt{a^2+b^2}=\sqrt{2c^2+c^2}=\sqrt{3}c\)
\(CB_1=\sqrt{b^2+c^2}=\sqrt{2}c\)
所以 \(AC\ne CB_1\)，\(C\) 错误.

对于 \(D\)，\(B_1D\) 与平面 \(BB_1C_1C\) 所成角为 \(\angle DB_1C\)，且
\(\sin\angle DB_1C=\dfrac{CD}{B_1D}=\dfrac{a}{2c}=\dfrac{\sqrt{2}c}{2c}=\dfrac{\sqrt{2}}{2}\)
又 \(0<\angle DB_1C<90^\circ\)，所以 \(\angle DB_1C=45^\circ\).
故 \(D\) 正确.
故选：\(D\).
\end{solution}

\begin{problem}
甲，乙两个圆锥的母线长相等，侧面展开图的圆心角之和为 \(2\pi\)，侧面积分别为 \(S_{\text{甲}}\) 和 \(S_{\text{乙}}\)，体积分别为 \(V_{\text{甲}}\) 和 \(V_{\text{乙}}\). 若 \(\dfrac{S_{\text{甲}}}{S_{\text{乙}}}=2\)，则 \(\dfrac{V_{\text{甲}}}{V_{\text{乙}}}=\)（\quad）

\choices
    {\(\sqrt{5}\)}
    {\(2\sqrt{2}\)}
    {\(\sqrt{10}\)}
    {\(\dfrac{5\sqrt{10}}{4}\)}

\end{problem}

\begin{answer}
C
\end{answer}

\begin{solution}
设母线长为 \(l\)，甲圆锥底面半径为 \(r_1\)，乙圆锥底面半径为 \(r_2\).
由圆锥的侧面积公式可得
\(\dfrac{S_{\text{甲}}}{S_{\text{乙}}} =\dfrac{\pi r_1l}{\pi r_2l} =\dfrac{r_1}{r_2} =2\)
所以
\(r_1=2r_2\).
又因为两个圆锥侧面展开图的圆心角之和为 \(2\pi\)，所以
\(\dfrac{2\pi r_1}{l}+\dfrac{2\pi r_2}{l}=2\pi\)
即
\(\dfrac{r_1+r_2}{l}=1\).
联立 \(r_1=2r_2\)，得
\(r_1=\dfrac{2}{3}l\)，\(r_2=\dfrac{1}{3}l\).
所以甲圆锥的高
\(h_1=\sqrt{l^2-r_1^2}=\sqrt{l^2-\dfrac{4}{9}l^2}=\dfrac{\sqrt{5}}{3}l\)
乙圆锥的高
\(h_2=\sqrt{l^2-r_2^2}=\sqrt{l^2-\dfrac{1}{9}l^2}=\dfrac{2\sqrt{2}}{3}l\).
于是 \(\dfrac{V_{\text{甲}}}{V_{\text{乙}}}=\dfrac{\dfrac{1}{3}\pi r_1^2h_1}{\dfrac{1}{3}\pi r_2^2h_2}=\dfrac{\dfrac{4}{9}l^2\times\dfrac{\sqrt{5}}{3}l}{\dfrac{1}{9}l^2\times\dfrac{2\sqrt{2}}{3}l}=\sqrt{10}\).
故选：\(C\).
\end{solution}

\begin{problem}
已知椭圆 \(C:\dfrac{x^2}{a^2}+\dfrac{y^2}{b^2}=1\)（\(a>b>0\)）的离心率为 \(\dfrac{1}{3}\)，\(A_1,A_2\) 分别为 \(C\) 的左，右顶点，\(B\) 为 \(C\) 的上顶点. 若 \(\overrightarrow{BA_1}\cdot\overrightarrow{BA_2}=-1\)，则 \(C\) 的方程为（\quad）

\choices
    {\(\dfrac{x^2}{18}+\dfrac{y^2}{16}=1\)}
    {\(\dfrac{x^2}{9}+\dfrac{y^2}{8}=1\)}
    {\(\dfrac{x^2}{3}+\dfrac{y^2}{2}=1\)}
    {\(\dfrac{x^2}{2}+y^2=1\)}

\end{problem}

\begin{answer}
B
\end{answer}

\begin{solution}
因为椭圆的离心率
\(e=\dfrac{c}{a}=\dfrac{1}{3}\)
所以
\(\dfrac{b^2}{a^2}=1-\dfrac{c^2}{a^2}=1-\dfrac{1}{9}=\dfrac{8}{9}\)
即
\(b^2=\dfrac{8}{9}a^2\).
又 \(A_1,A_2\) 分别为椭圆 \(C\) 的左，右顶点，\(B\) 为上顶点，所以
\(A_1\left(-a,0\right)\)，\(A_2\left(a,0\right)\)，\(B\left(0,b\right)\).
于是
\(\overrightarrow{BA_1}=\left(-a,-b\right)\)，\(\overrightarrow{BA_2}=\left(a,-b\right)\).
由 \(\overrightarrow{BA_1}\cdot\overrightarrow{BA_2}=-1\) 得
\(-a^2+b^2=-1\).
将 \(b^2=\dfrac{8}{9}a^2\) 代入，得
\(-a^2+\dfrac{8}{9}a^2=-1\)
解得
\(a^2=9\)，\(b^2=8\).
所以椭圆 \(C\) 的方程为
\(\dfrac{x^2}{9}+\dfrac{y^2}{8}=1\).
故选：\(B\).
\end{solution}

\begin{problem}
已知 \(9^m=10\)，\(a=10^m-11\)，\(b=8^m-9\)，则（\quad）

\choices
    {\(a>0>b\)}
    {\(a>b>0\)}
    {\(b>a>0\)}
    {\(b>0>a\)}

\end{problem}

\begin{answer}
A
\end{answer}

\begin{solution}
由 \(9^m=10\) 可得
\(m=\log_910=\dfrac{\lg10}{\lg9}>1\).
又 \(\lg9\lg11<\left(\dfrac{\lg9+\lg11}{2}\right)^2=\left(\dfrac{\lg99}{2}\right)^2<1=\left(\lg10\right)^2\)
所以
\(\dfrac{\lg10}{\lg9}>\dfrac{\lg11}{\lg10}\)
即
\(m>\lg11\).
于是
\(a=10^m-11>10^{\lg11}-11=0\).
再有 \(\lg8\lg10<\left(\dfrac{\lg8+\lg10}{2}\right)^2=\left(\dfrac{\lg80}{2}\right)^2<\left(\lg9\right)^2\)
所以
\(\dfrac{\lg10}{\lg9}<\dfrac{\lg9}{\lg8}=\log_89\)
即
\(m<\log_89\).
于是
\(b=8^m-9<8^{\log_89}-9=0\).
综上，
\(a>0>b\).
故选：\(A\).
\end{solution}
\section{填空题}

\begin{problem}
已知向量 \(\bs a=\left(m,3\right)\)，\(\bs b=\left(1,m+1\right)\). 若 \(\bs a\perp\bs b\)，则 \(m=\)\fillinblank{}

\end{problem}

\begin{answer}
\(-\dfrac{3}{4}\)
\end{answer}

\begin{solution}
因为 \(\bs a\perp\bs b\)，所以
\(\bs a\cdot\bs b=m+3\left(m+1\right)=0\)
即
\(4m+3=0\)
解得
\(m=-\dfrac{3}{4}\).
故答案为：\(-\dfrac{3}{4}\).
\end{solution}

\begin{problem}
设圆 \(M\) 的圆心在直线 \(2x+y-1=0\) 上，点 \(\left(3,0\right)\) 和 \(\left(0,1\right)\) 均在圆 \(M\) 上，则圆 \(M\) 的方程为\fillinblank{}.

\end{problem}

\begin{answer}
\(\left(x-1\right)^2+\left(y+1\right)^2=5\)
\end{answer}

\begin{solution}
设圆心 \(M\) 的坐标为 \(\left(a,1-2a\right)\).
又因为点 \(\left(3,0\right)\) 和 \(\left(0,1\right)\) 均在圆 \(M\) 上，所以圆心到这两点的距离相等，即
\(\left(a-3\right)^2+\left(1-2a\right)^2=a^2+\left(-2a\right)^2\).
化简得
\(a=1\).
所以圆心为
\(M\left(1,-1\right)\).
半径满足
\(R^2=\left(1-3\right)^2+\left(-1\right)^2=5\).
故圆 \(M\) 的方程为
\(\left(x-1\right)^2+\left(y+1\right)^2=5\).
故答案为：\(\left(x-1\right)^2+\left(y+1\right)^2=5\).
\end{solution}

\begin{problem}
记双曲线 \(C:\dfrac{x^2}{a^2}-\dfrac{y^2}{b^2}=1\)（\(a>0\)，\(b>0\)）的离心率为 \(e\)，写出满足条件“直线 \(y=2x\) 与 \(C\) 无公共点”的 \(e\) 的一个值\fillinblank{}.

\end{problem}

\begin{answer}
2（满足 \(1<e\le \sqrt{5}\) 皆可）
\end{answer}

\begin{solution}
双曲线
\(C:\dfrac{x^2}{a^2}-\dfrac{y^2}{b^2}=1\)
的渐近线方程为
\(y=\pm\dfrac{b}{a}x\).
要使直线 \(y=2x\) 与双曲线 \(C\) 无公共点，只需
\(0<\dfrac{b}{a}\le2\)
即
\(\dfrac{b^2}{a^2}\le4\).
又
\(e=\dfrac{c}{a}=\sqrt{1+\dfrac{b^2}{a^2}}\)
所以
\(1<e\le\sqrt{5}\).
因为 \(2\) 满足 \(1<e\le\sqrt{5}\)，所以 \(e=2\) 符合题意.
故答案为：2（满足 \(1<e\le \sqrt{5}\) 皆可）.
\end{solution}

\begin{problem}
已知 \(\triangle ABC\) 中，点 \(D\) 在边 \(BC\) 上，\(\angle ADB=120^\circ\)，\(AD=2\)，且 \(CD=2BD\). 当 \(\dfrac{AC}{AB}\) 取得最小值时，\(BD=\)\fillinblank{}.

\end{problem}

\begin{answer}
\(\sqrt{3}-1\)
\end{answer}

\begin{solution}
设 \(CD=2BD=2m>0\)，则 \(BD=m\).
在 \(\triangle ABD\) 中，由余弦定理得
\(AB^2=BD^2+AD^2-2BD\cdot AD\cos\angle ADB =m^2+4+2m\).
在 \(\triangle ACD\) 中，由余弦定理得
\(AC^2=CD^2+AD^2-2CD\cdot AD\cos\angle ADC =4m^2+4-4m\).
所以 \(\dfrac{AC^2}{AB^2}=\dfrac{4m^2+4-4m}{m^2+4+2m}=4-\dfrac{12\left(m+1\right)}{\left(m+1\right)^2+3}\).
又
\(\left(m+1\right)^2+3\ge2\sqrt{3}\left(m+1\right)\)
所以 \(\dfrac{12\left(m+1\right)}{\left(m+1\right)^2+3}\le\dfrac{12\left(m+1\right)}{2\sqrt{3}\left(m+1\right)}=2\sqrt{3}\).
因此
\(\dfrac{AC^2}{AB^2}\ge4-2\sqrt{3}\).
当且仅当
\(m+1=\dfrac{\sqrt{3}}{m+1}\)
即
\(m=\sqrt{3}-1\)
时等号成立.
所以当 \(\dfrac{AC}{AB}\) 取得最小值时，
\(BD=m=\sqrt{3}-1\).
故答案为：\(\sqrt{3}-1\).
\end{solution}
\section{解答题}

\begin{problem}
甲，乙两城之间的长途客车均由 \(A\) 和 \(B\) 两家公司运营，为了解这两家公司长途客车的运行情况，随机调查了甲，乙两城之间的 \(500\) 个班次，得到下面列联表：

\begin{center}
\begin{tabular}{c c c}
\hline
& 准点班次数 & 未准点班次数 \\
\hline
\(A\) & 240 & 20 \\
\(B\) & 210 & 30 \\
\hline
\end{tabular}
\end{center}

\begin{enumerate}
\item 根据上表，分别估计这两家公司甲，乙两城之间的长途客车准点的概率；
\item 能否有 \(90\%\) 的把握认为甲，乙两城之间的长途客车是否准点与客车所属公司有关？
\end{enumerate}

附：\(K^2=\dfrac{n\left(ad-bc\right)^2}{\left(a+b\right)\left(c+d\right)\left(a+c\right)\left(b+d\right)}\)

\begin{center}
\begin{tabular}{c c c c}
\hline
\(P\left(K^2\ge k\right)\) & 0.100 & 0.050 & 0.010 \\
\hline
\(k\) & 2.706 & 3.841 & 6.635 \\
\hline
\end{tabular}
\end{center}

\end{problem}

\begin{solution}
（1）根据表中数据，\(A\) 公司共有班次 \(260\) 次，准点班次有 \(240\) 次，
设 \(A\) 家公司长途客车准点事件为 \(M\)，则
\(P\left(M\right)=\dfrac{240}{260}=\dfrac{12}{13}\).
\(B\) 公司共有班次 \(240\) 次，准点班次有 \(210\) 次，
设 \(B\) 家公司长途客车准点事件为 \(N\)，则
\(P\left(N\right)=\dfrac{210}{240}=\dfrac{7}{8}\).
所以 \(A\) 家公司长途客车准点的概率为 \(\dfrac{12}{13}\)，\(B\) 家公司长途客车准点的概率为 \(\dfrac{7}{8}\).

（2）根据题意列联表为

\begin{center}
\begin{tabular}{c c c c}
\hline
& 准点班次数 & 未准点班次数 & 合计 \\
\hline
\(A\) & 240 & 20 & 260 \\
\(B\) & 210 & 30 & 240 \\
\hline
合计 & 450 & 50 & 500 \\
\hline
\end{tabular}
\end{center}

所以 \(K^2=\dfrac{500\left(240\times30-210\times20\right)^2}{260\times240\times450\times50}\approx3.205\).
因为
\(3.205>2.706\)
所以根据临界值表可知，有 \(90\%\) 的把握认为甲，乙两城之间的长途客车是否准点与客车所属公司有关.
\end{solution}

\begin{problem}
记 \(S_n\) 为数列 \(\{a_n\}\) 的前 \(n\) 项和. 已知 \(\dfrac{2S_n}{n}+n=2a_n+1\).

\begin{enumerate}
\item 证明：\(\{a_n\}\) 是等差数列；
\item 若 \(a_4,a_7,a_9\) 成等比数列，求 \(S_n\) 的最小值.
\end{enumerate}

\end{problem}

\begin{solution}
（1）因为
\(\dfrac{2S_n}{n}+n=2a_n+1\)
所以
\(2S_n+n^2=2na_n+n \circled{1}\).
当 \(n\ge2\) 时，
\(2S_{n-1}+\left(n-1\right)^2=2\left(n-1\right)a_{n-1}+\left(n-1\right) \circled{2}\).
\(\circled{1}-\circled{2}\) 得
\(2S_n+n^2-2S_{n-1}-\left(n-1\right)^2 =2na_n+n-2\left(n-1\right)a_{n-1}-\left(n-1\right)\).
又因为
\(S_n-S_{n-1}=a_n\)
所以
\(2a_n+2n-1=2na_n-2\left(n-1\right)a_{n-1}+1\).
化简得
\(2\left(n-1\right)a_n-2\left(n-1\right)a_{n-1}=2\left(n-1\right)\).
因此
\(a_n-a_{n-1}=1\)，\(n\ge2\).
故 \(\{a_n\}\) 是以 \(1\) 为公差的等差数列.

（2）由（1）可得
\(a_4=a_1+3\)，\(a_7=a_1+6\)，\(a_9=a_1+8\).
又因为 \(a_4,a_7,a_9\) 成等比数列，所以
\(a_7^2=a_4a_9\)
即
\(\left(a_1+6\right)^2=\left(a_1+3\right)\left(a_1+8\right)\).
解得
\(a_1=-12\).
所以
\(a_n=n-13\).
从而 \(S_n=\sum_{k=1}^{n}\left(k-13\right)=-12n+\dfrac{n\left(n-1\right)}{2}=\dfrac{1}{2}n^2-\dfrac{25}{2}n\).
于是
\(S_n=\dfrac{1}{2}\left(n-\dfrac{25}{2}\right)^2-\dfrac{625}{8}\).
所以当 \(n=12\) 或 \(n=13\) 时，\(S_n\) 取最小值
\(-78\).
故 \(S_n\) 的最小值为 \(-78\).
\end{solution}

\begin{problem}
小明同学参加综合实践活动，设计了一个封闭的包装盒，包装盒如图所示：底面 \(ABCD\) 是边长为 \(8\)（单位：cm）的正方形，\(EAB,FBC,GCD,HDA\) 均为正三角形，且它们所在的平面都与平面 \(ABCD\) 垂直.

\begin{enumerate}
\item 证明：\(EF//\) 平面 \(ABCD\)；
\item 求该包装盒的容积（不计包装盒材料的厚度）.
\end{enumerate}

\bitmapfigure[max width=6.0cm,max totalheight=4.8cm]{img/2022/national_paper_a_liberal/q19_fig1.png}

\end{problem}

\begin{solution}
（1）分别取 \(AB,BC\) 的中点 \(M,N\)，连接 \(MN\). 如图所示：

\bitmapfigure[max width=5.0cm,max totalheight=4.8cm]{img/2022/national_paper_a_liberal/q19_solution_fig1.png}

因为 \(EAB,FBC\) 为全等的正三角形，所以
\(EM\perp AB\)，\(FN\perp BC\)，\(EM=FN\).
又因为平面 \(EAB\perp\) 平面 \(ABCD\)，平面 \(EAB\cap\) 平面 \(ABCD=AB\)，且 \(EM\subset\) 平面 \(EAB\)，
所以
\(EM\perp \text{平面 }ABCD\).
再对侧面 \(FBC\) 作相同的中点高线分析，可得
\(FN\perp \text{平面 }ABCD\).
根据线面垂直的性质定理可知
\(EM//FN\).
又 \(EM=FN\)，所以四边形 \(EMNF\) 为平行四边形，
故
\(EF//MN\).
又 \(EF\not\subset\) 平面 \(ABCD\)，\(MN\subset\) 平面 \(ABCD\)，
所以
\(EF//\text{平面 }ABCD\).

（2）分别取 \(AD,DC\) 的中点 \(K,L\). 如图所示：

\bitmapfigure[max width=5.0cm,max totalheight=4.8cm]{img/2022/national_paper_a_liberal/q19_solution_fig2.png}

由（1）知，\(EF//MN\) 且 \(EF=MN\). 由四个侧面均为边长为 \(8\) 的正三角形，且对应中点构造完全一致，可得
\(HE//KM,\quad HE=KM\)，
\(HG//KL,\quad HG=KL\)，
\(GF//LN,\quad GF=LN\).
由平面知识可知
\(BD\perp MN\)，\(MN\perp MK\)，
\(KM=MN=NL=LK\).
所以该几何体的体积等于长方体 \(KMNL-EFGH\) 的体积加上四棱锥 \(B-MNFE\) 体积的 \(4\) 倍.
因为
\(MN=NL=LK=KM=4\sqrt{2}\)，
\(EM=8\sin60^\circ=4\sqrt{3}\)，
点 \(B\) 到平面 \(MNFE\) 的距离即为点 \(B\) 到直线 \(MN\) 的距离，
\(d=2\sqrt{2}\).
所以该几何体的体积 \(V=\left(4\sqrt{2}\right)^2\times4\sqrt{3}+4\times\dfrac{1}{3}\times4\sqrt{2}\times4\sqrt{3}\times2\sqrt{2}=128\sqrt{3}+\dfrac{256\sqrt{3}}{3}=\dfrac{640\sqrt{3}}{3}\).
所以该包装盒的容积为 \(\dfrac{640\sqrt{3}}{3}\).
\end{solution}

\begin{problem}
已知函数 \(f(x)=x^3-x\)，\(g(x)=x^2+a\)，曲线 \(y=f(x)\) 在点 \(\left(x_1,f(x_1)\right)\) 处的切线也是曲线 \(y=g(x)\) 的切线.

\begin{enumerate}
\item 若 \(x_1=-1\)，求 \(a\)；
\item 求 \(a\) 的取值范围.
\end{enumerate}

\end{problem}

\begin{solution}
（1）由题意知
\(f(-1)=\left(-1\right)^3-\left(-1\right)=0\)，
\(f'(x)=3x^2-1\)，
\(f'(-1)=2\).
所以曲线 \(y=f(x)\) 在点 \(\left(-1,0\right)\) 处的切线方程为
\(y=2\left(x+1\right)\)
即
\(y=2x+2\).
设该切线与 \(g(x)\) 切于点 \(\left(x_2,g(x_2)\right)\)，
因为
\(g'(x)=2x\)
所以
\(g'(x_2)=2x_2=2\)
解得
\(x_2=1\).
又
\(g(1)=1+a=2+2\)
解得
\(a=3\).

（2）由
\(f'(x)=3x^2-1\)
可知曲线 \(y=f(x)\) 在点 \(\left(x_1,f(x_1)\right)\) 处的切线方程为
\(y-\left(x_1^3-x_1\right)=\left(3x_1^2-1\right)\left(x-x_1\right)\)
整理得
\(y=\left(3x_1^2-1\right)x-2x_1^3\).
设该切线与 \(g(x)\) 切于点 \(\left(x_2,g(x_2)\right)\)，则 \(g'(x_2)=2x_2\)，其切线方程为
\(y-\left(x_2^2+a\right)=2x_2\left(x-x_2\right)\)
整理得
\(y=2x_2x-x_2^2+a\).
因为两条切线重合，所以 \(3x_1^2-1=2x_2\)，\(-2x_1^3=-x_2^2+a\).
由第一式得
\(x_2=\dfrac{3x_1^2-1}{2}\).
代入第二式得
\(a=x_2^2-2x_1^3 =\left(\dfrac{3x_1^2-1}{2}\right)^2-2x_1^3\)
\(=\dfrac{9}{4}x_1^4-2x_1^3-\dfrac{3}{2}x_1^2+\dfrac{1}{4}\).
令
\(h(x)=\dfrac{9}{4}x^4-2x^3-\dfrac{3}{2}x^2+\dfrac{1}{4}\)，
则
\(h'(x)=9x^3-6x^2-3x=3x\left(3x+1\right)\left(x-1\right)\).
所以当 \(x<-\dfrac{1}{3}\) 时，\(h'(x)<0\)；

当 \(-\dfrac{1}{3}<x<0\) 时，\(h'(x)>0\)；

当 \(0<x<1\) 时，\(h'(x)<0\)；

当 \(x>1\) 时，\(h'(x)>0\).
故 \(h(x)\) 的最小值在 \(x=1\) 处取得，且
\(h(1)=\dfrac{9}{4}-2-\dfrac{3}{2}+\dfrac{1}{4}=-1\).
所以
\(a\ge-1\).
故 \(a\) 的取值范围为 \(\left[-1,+\infty\right)\).
\end{solution}

\begin{problem}
设抛物线 \(C:y^2=2px\)（\(p>0\)）的焦点为 \(F\)，点 \(D\left(p,0\right)\)，过 \(F\) 的直线交 \(C\) 于 \(M,N\) 两点. 当直线 \(MD\) 垂直于 \(x\) 轴时，\(MF=3\).

\begin{enumerate}
\item 求 \(C\) 的方程；
\item 设直线 \(MD,ND\) 与 \(C\) 的另一个交点分别为 \(A,B\)，记直线 \(MN,AB\) 的倾斜角分别为 \(\alpha,\beta\). 当 \(\alpha-\beta\) 取得最大值时，求直线 \(AB\) 的方程.
\end{enumerate}

\end{problem}

\begin{solution}
（1）抛物线 \(C:y^2=2px\) 的准线为
\(x=-\dfrac{p}{2}\).
当 \(MD\perp x\) 轴时，点 \(M\) 的横坐标为 \(p\).
由抛物线的定义可知
\(MF=p+\dfrac{p}{2}=3\).
所以
\(p=2\).
故抛物线 \(C\) 的方程为
\(y^2=4x\).

（2）设 \(M\left(\dfrac{y_1^2}{4},y_1\right)\)，\(N\left(\dfrac{y_2^2}{4},y_2\right)\)，\(A\left(\dfrac{y_3^2}{4},y_3\right)\)，\(B\left(\dfrac{y_4^2}{4},y_4\right)\).
记
\(x_1=\dfrac{y_1^2}{4}\)，\(x_2=\dfrac{y_2^2}{4}\).
设直线 \(MN\) 的方程为
\(x=my+1\).
与抛物线 \(y^2=4x\) 联立，得
\(y^2-4my-4=0\).
所以
\(y_1y_2=-4\).
由斜率公式可得
\(k_{MN}=\dfrac{y_1-y_2}{\dfrac{y_1^2}{4}-\dfrac{y_2^2}{4}}=\dfrac{4}{y_1+y_2}\).
又因为直线 \(MD\) 经过点 \(D\left(2,0\right)\) 和点 \(M\left(\dfrac{y_1^2}{4},y_1\right)\)，所以
\(x=\dfrac{x_1-2}{y_1}y+2\).
将其代入 \(y^2=4x\)，得
\(y^2-\dfrac{4\left(x_1-2\right)}{y_1}y-8=0\).
所以
\(y_1y_3=-8\)
从而
\(y_3=2y_2\).
又因为直线 \(ND\) 经过点 \(D\left(2,0\right)\) 和点 \(N\left(\dfrac{y_2^2}{4},y_2\right)\)，所以
\(x=\dfrac{x_2-2}{y_2}y+2\).
将其代入 \(y^2=4x\)，得
\(y^2-\dfrac{4\left(x_2-2\right)}{y_2}y-8=0\).
所以
\(y_2y_4=-8\)
从而
\(y_4=2y_1\).
于是 \(k_{AB}=\dfrac{y_3-y_4}{\dfrac{y_3^2}{4}-\dfrac{y_4^2}{4}}=\dfrac{4}{y_3+y_4}=\dfrac{4}{2\left(y_1+y_2\right)}=\dfrac{k_{MN}}{2}\).
因为直线 \(MN,AB\) 的倾斜角分别为 \(\alpha,\beta\)，所以
\(\tan\alpha=k_{MN}\)，\(\tan\beta=k_{AB}=\dfrac{k_{MN}}{2}\).
设
\(k_{AB}=k>0\)，
则
\(k_{MN}=2k\).
于是 \(\tan\left(\alpha-\beta\right)=\dfrac{\tan\alpha-\tan\beta}{1+\tan\alpha\tan\beta}=\dfrac{2k-k}{1+2k^2}=\dfrac{k}{1+2k^2}\).
又 \(\dfrac{k}{1+2k^2}=\dfrac{1}{\dfrac{1}{k}+2k}\le\dfrac{1}{2\sqrt{\dfrac{1}{k}\cdot2k}}=\dfrac{\sqrt{2}}{4}\).
当且仅当
\(\dfrac{1}{k}=2k\)
即
\(k=\dfrac{\sqrt{2}}{2}\)
时等号成立.
所以当 \(\alpha-\beta\) 取得最大值时，
\(k_{AB}=\dfrac{\sqrt{2}}{2}\).
设直线 \(AB\) 的方程为
\(x=\dfrac{\sqrt{2}}{2}y+n\).
代入抛物线方程 \(y^2=4x\)，得
\(y^2-2\sqrt{2}y-4n=0\).
于是
\(y_3y_4=-4n\).
又由前面已知
\(y_3y_4=4y_1y_2=4\times\left(-4\right)=-16\)
所以
\(n=4\).
故直线 \(AB\) 的方程为
\(x=\dfrac{\sqrt{2}}{2}y+4\).
\end{solution}

\begin{problem}
在直角坐标系 \(xOy\) 中，曲线 \(C_1\) 的参数方程为 \(x=\dfrac{2+t^2}{6}\)，\(y=t\)（\(t\) 为参数）；曲线 \(C_2\) 的参数方程为 \(x=-\dfrac{2+s^2}{6}\)，\(y=-s\)（\(s\) 为参数）.

\begin{enumerate}
\item 写出 \(C_1\) 的普通方程；
\item 以坐标原点为极点，\(x\) 轴正半轴为极轴建立极坐标系，曲线 \(C_3\) 的极坐标方程为 \(2\cos\theta-\sin\theta=0\)，求 \(C_3\) 与 \(C_1\) 交点的直角坐标，及 \(C_3\) 与 \(C_2\) 交点的直角坐标.
\end{enumerate}

\end{problem}

\begin{solution}
（1）因为
\(x=\dfrac{2+t^2}{6}\)，\(y=t\)，
所以
\(x=\dfrac{2+y^2}{6}\).
故 \(C_1\) 的普通方程为
\(y^2=6x-2\)（\(y\ge0\)）.

（2）由 \(x=-\dfrac{2+s^2}{6}\)，\(y=-s\) 得
\(6x=-2-y^2\)，
所以 \(C_2\) 的普通方程为
\(y^2=-6x-2\)（\(y\le0\)）.
又因为
\(2\cos\theta-\sin\theta=0\)
所以
\(2\rho\cos\theta-\rho\sin\theta=0\)，
即 \(C_3\) 的普通方程为
\(2x-y=0\).
联立 \(y^2=6x-2\) 与 \(2x-y=0\)，即 \(y^2=3y-2\)，\(x=\dfrac{y}{2}\).
解得
\(\left(\dfrac{1}{2},1\right)\)，\(\left(1,2\right)\).
所以 \(C_3\) 与 \(C_1\) 的交点坐标为
\(\left(\dfrac{1}{2},1\right)\)，\(\left(1,2\right)\).
再联立 \(y^2=-6x-2\) 与 \(2x-y=0\)，即 \(y^2=-3y-2\)，\(x=\dfrac{y}{2}\).
解得
\(\left(-\dfrac{1}{2},-1\right)\)，\(\left(-1,-2\right)\).
所以 \(C_3\) 与 \(C_2\) 的交点坐标为
\(\left(-\dfrac{1}{2},-1\right)\)，\(\left(-1,-2\right)\).
\end{solution}

\begin{problem}
已知 \(a,b,c\) 均为正数，且 \(a^2+b^2+4c^2=3\)，证明：

\begin{enumerate}
\item \(a+b+2c\le3\)；
\item 若 \(b=2c\)，则 \(\dfrac{1}{a}+\dfrac{1}{c}\ge3\).
\end{enumerate}

\end{problem}

\begin{solution}
（1）由柯西不等式有
\(\left[a^2+b^2+\left(2c\right)^2\right]\left(1^2+1^2+1^2\right)\ge\left(a+b+2c\right)^2\).
又因为
\(a^2+b^2+4c^2=3\)，
所以
\(\left(a+b+2c\right)^2\le9\).
因为 \(a,b,c>0\)，所以
\(a+b+2c\le3\).

（2）因为 \(b=2c\)，且 \(a>0\)，\(b>0\)，\(c>0\)，由（1）得
\(a+b+2c=a+4c\le3\).
所以
\(0<a+4c\le3\)，
从而
\(\dfrac{1}{a+4c}\ge\dfrac{1}{3}\).
又由权方和不等式知 \(\dfrac{1}{a}+\dfrac{1}{c}=\dfrac{1^2}{a}+\dfrac{2^2}{4c}\ge\dfrac{\left(1+2\right)^2}{a+4c}=\dfrac{9}{a+4c}\ge3\).
故
\(\dfrac{1}{a}+\dfrac{1}{c}\ge3\).
\end{solution}
